\section[Confiabilidad dinámica y niveles de evidencia]
{Confiabilidad dinámica: predictibilidad, caos transitorio y
\mbox{niveles} de evidencia}
\label{sec:confiabilidad-dinamica}

\HAFOCardReliability

La confiabilidad dinámica se descompone en cuatro problemas: si la trayectoria
calculada aproxima el PVI declarado; si el comportamiento irregular persiste al ampliar
la ventana; si existe un conjunto invariante que atrae los datos declarados; y
si su cuenca queda separada de una vecindad de cada equilibrio. Para ODE la
atracción local puede formularse sobre un abierto del espacio de estados; en
Caputo se especifica la familia de reinicios o un espacio admisible invariante
\(X_{\mathrm{adm}}\) con su topología relativa. Cada propiedad requiere criterios propios, desde el control del error hasta
la separación de cuencas.

La teoría de Lyapunov permite estudiar la predictibilidad y la expansión de
perturbaciones \cite{VallejoSanjuan2017Predictability,DatserisParlitz2022Nonlinear}.
Los tiempos de escape y las crisis distinguen la recurrencia transitoria de
la atracción \cite{Ott2002Chaos,TelGruiz2006Chaotic}. En sistemas
fraccionarios, ambos análisis deben incorporar el operador y la memoria del
problema de Caputo \cite{SunHeWang2022FractionalChaos,Sabatier2021InfiniteMemory}.

\subsection{Error de trayectoria y horizonte de predictibilidad}

Considérese primero la EDO autónoma
\begin{equation}
 \dot x=f(x),\qquad x(t_0)=x_0,
 \label{eq:reliability-ode}
\end{equation}
con \(f\) Lipschitz de constante \(L\) en la región recorrida. Una curva
numérica continua y diferenciable por tramos \(\widetilde x\) posee el defecto
\begin{equation}
 r_h(t)=\dot{\widetilde x}(t)-f(\widetilde x(t)).
 \label{eq:reliability-defect}
\end{equation}
Si \(e=x-\widetilde x\), integrar la diferencia de las ecuaciones da
\[
 e(t)=e(t_0)+\int_{t_0}^t[f(x(s))-f(\widetilde x(s))]\,\dd s
                  -\int_{t_0}^tr_h(s)\,\dd s.
\]
Por Lipschitz y la desigualdad triangular,
\[
 u(t):=\|e(t)\|\leq\|e(t_0)\|
            +\int_{t_0}^t[Lu(s)+\|r_h(s)\|]\,\dd s=:v(t).
\]
La función \(v\) es absolutamente continua y cumple
\(v'\leq Lv+\|r_h\|\) casi en todas partes. Multiplicar por
\(\e^{-L(t-t_0)}\) e integrar conduce a
\begin{equation}
 \|e(t)\|
 \leq \e^{L(t-t_0)}\|e(t_0)\|
 +\int_{t_0}^{t}\e^{L(t-s)}\|r_h(s)\|\,\dd s .
 \label{eq:reliability-gronwall}
\end{equation}
La convergencia local del integrador no impide que el error de una trayectoria
individual crezca rápidamente. En un régimen con exponente principal
\(\lambda_1>0\), la estimación heurística
\begin{equation}
 T_{\mathrm{pred}}
 \approx\frac{1}{\lambda_1}
 \log\!\left(\frac{\varepsilon_{\mathrm{adm}}}
                   {\varepsilon_0}\right)
 \label{eq:reliability-predictability-horizon}
\end{equation}
se obtiene de \(\varepsilon_{\mathrm{adm}}\approx
\varepsilon_0\e^{\lambda_1T_{\mathrm{pred}}}\), al tomar logaritmos, con
\(0<\varepsilon_0<\varepsilon_{\mathrm{adm}}\). Relaciona incertidumbre
inicial, tolerancia admisible y horizonte de predictibilidad. No es un teorema universal: depende de la órbita,
la dirección de perturbación, la norma, la ventana y el régimen dinámico
\cite{VallejoSanjuan2017Predictability}.

Después de \(T_{\mathrm{pred}}\), la comparación puede centrarse en
estadísticas convergentes: tiempos de escape, clases terminales,
distribuciones de retorno, espectros y exponentes en varias ventanas. La
semejanza visual de las nubes requiere el contraste de esos observables.

\subsection{Alcance y limitaciones del sombreado}

Sea \(P:X\to X\) un mapa sobre un espacio métrico. Una sucesión
\((z_k)\) es una \(\delta\)-pseudoórbita si
\begin{equation}
 d(P(z_k),z_{k+1})\leq\delta .
 \label{eq:reliability-pseudo-orbit}
\end{equation}
Está \(\varepsilon\)-sombreada si existe \(y\in X\) tal que
\begin{equation}
 d(P^k(y),z_k)\leq\varepsilon
 \quad\text{para todos los índices de la ventana considerada}.
 \label{eq:reliability-shadowing}
\end{equation}
Los teoremas de \emph{shadowing} requieren hipótesis, típicamente
hiperbolicidad o variantes controladas. Cuando aplican, aseguran proximidad a
\emph{alguna} órbita exacta; no obligan a que esa órbita empiece en el dato
nominal ni a que pertenezca a la misma cuenca cerca de una frontera
\cite{Ott2002Chaos,VallejoSanjuan2017Predictability}.

La precisión nominal del integrador no sustituye un argumento de sombreado. La
evaluación registra paso, tolerancias, segundo método,
residuo cuando es accesible y estabilidad de los observables. Cerca de una
frontera de cuenca, una perturbación numérica pequeña puede seleccionar un
destino distinto. Su sensibilidad debe cuantificarse.

\subsection{Exponentes de Lyapunov de tiempo finito}

Para \cref{eq:reliability-ode}, sea \(Y\) la matriz fundamental
\begin{equation}
 \dot Y(t)=Df(x(t))Y(t),\qquad Y(t_0)=I.
 \label{eq:reliability-variational}
\end{equation}
Una perturbación infinitesimal inicial \(v\) evoluciona como \(Yv\).
El cociente de amplificación satisface
\[
 \frac{\|Yv\|_2^2}{\|v\|_2^2}
 =\frac{v^{\mathsf T}Y^{\mathsf T}Yv}{v^{\mathsf T}v}.
\]
Los extremos principales son los autovalores de la matriz simétrica positiva
\(Y^{\mathsf T}Y\); sus raíces cuadradas son los valores singulares
\(\sigma_i(Y)\). Tomar el logaritmo de cada factor de amplificación y
dividir por la duración define los exponentes de tiempo finito (FTLE):
\begin{equation}
 \Lambda_i(t_0,T;x_0)
 =\frac1T\log\sigma_i\!\left(Y(t_0+T;t_0,x_0)\right).
 \label{eq:reliability-ftle}
\end{equation}
A diferencia del límite asintótico de
\cref{geo:eq:lyapunov-definition}, \(\Lambda_i\) depende de \(x_0\), \(t_0\),
\(T\), la norma y el tratamiento numérico. Esta dependencia es informativa:
su distribución a lo largo de ventanas revela intermitencia, regiones
pegajosas y pérdida temporal de hiperbolicidad
\cite{VallejoSanjuan2017Predictability}.

Un FTLE principal positivo indica expansión durante la ventana considerada.
Se necesitan argumentos adicionales para establecer:
\begin{itemize}
 \item que el exponente asintótico exista o sea positivo;
 \item que la trayectoria pertenezca a un atractor;
 \item que haya herradura, entropía topológica positiva o transitividad;
 \item que la cuenca sea oculta.
\end{itemize}
La verificación mínima conserva la longitud y el solapamiento de las ventanas,
el intervalo de reortogonalización, la base tangente inicial, el transitorio
descartado y la convergencia al variar paso y horizonte. Una media sin su
distribución puede ocultar ventanas de escape.

\subsection{El defecto de Volterra en Caputo}

Para
\begin{equation}
 {}^{\mathrm C}D_{t_0}^{q}x(t)=f(x(t)),\qquad 0<q<1,
 \label{eq:reliability-caputo}
\end{equation}
el defecto apropiado a la forma de Volterra es
\begin{equation}
 \rho_h(t)=\widetilde x(t)-x_0
 -\frac1{\Gamma(q)}
 \int_{t_0}^{t}(t-s)^{q-1}f(\widetilde x(s))\,\dd s .
 \label{eq:reliability-volterra-defect}
\end{equation}
Restar esta identidad de la ecuación exacta de Volterra produce
\[
 x(t)-\widetilde x(t)=-\rho_h(t)+\frac1{\Gamma(q)}
 \int_{t_0}^t(t-s)^{q-1}[f(x(s))-f(\widetilde x(s))]\,\dd s.
\]
Si el PVI tiene solución en el intervalo considerado y \(f\) es Lipschitz
de constante \(L\) en una región que contiene ambas curvas, su norma satisface
\(u\leq\|\rho_h\|_\infty+LI^qu\). La iteración de integrales fraccionarias
empleada en \cref{camp:prop-soe-error-finite}, trasladando el origen a \(t_0\),
suma la serie de Mittag--Leffler y da
\cite{Lin2013GeneralizedGronwall}
\begin{equation}
 \|x(t)-\widetilde x(t)\|
 \leq
 \|\rho_h\|_{\infty,[t_0,t]}
 E_q\!\left(L(t-t_0)^q\right).
 \label{eq:reliability-fractional-gronwall}
\end{equation}
La cota contiene la función de Mittag--Leffler y el defecto acumulado desde
\(t_0\), incluidas las contribuciones anteriores a una ventana tardía.

Los modelos fraccionarios admiten realizaciones con un continuo de escalas y
una memoria efectivamente infinita \cite{Sabatier2021InfiniteMemory}. Por eso,
comparar ABM--PECE, aproximación frecuencial o Adomian es útil para estudiar
sesgo algorítmico \cite{SunHeWang2022FractionalChaos}, pero sólo son
comparables si resuelven el mismo operador, el mismo terminal inferior y la
misma inicialización. Un truncamiento de memoria define otro experimento salvo
que su error esté controlado.

La ecuación variacional fraccionaria es la de
\cref{eq:geofo-fractional-variational}; también conserva historia. Reiniciar
vectores tangentes como en una ODE ordinaria sin especificar cómo se transporta
su memoria no reproduce automáticamente esa ecuación. El cálculo de exponentes fraccionarios puede contrastarse con
Danca--Kuznetsov \cite{DancaKuznetsov2018Lyapunov}. Las secciones proyectadas
proporcionan detectores numéricos de retorno
\cite{ClementeLopezEtAl2022Poincare}; su formulación dinámica rigurosa
pertenece al espacio de historias.

\subsection{Silla caótica, escape y supervivencia}

Sea \(D\) una región compacta de un flujo entero \(\varphi_t\). El conjunto
invariante maximal en \(D\) es
\begin{equation}
 \operatorname{Inv}(D)
 =\{x\in D:\varphi_t(x)\in D\ \text{para todo }t\in\mathbb R\}.
 \label{eq:reliability-maximal-invariant}
\end{equation}
Una \emph{silla caótica} es, en términos operativos, un conjunto caótico
invariante no atractor con direcciones estables e inestables. Las condiciones
iniciales próximas a su conjunto estable pueden permanecer mucho tiempo en
la región, mostrar expansión y recurrencia, y después escapar a un equilibrio,
un ciclo, otro atractor o infinito \cite{Ott2002Chaos,TelGruiz2006Chaotic}.

Para una condición inicial \(x_0\in D\), se define el tiempo de primera salida
\begin{equation}
 \tau_D(x_0)=
 \inf\{t\geq0:\varphi_t(x_0)\notin D\},
 \label{eq:reliability-exit-time}
\end{equation}
con la convención \(\inf\varnothing=\infty\): este valor corresponde a una
órbita exacta que nunca sale, no a una integración finita sin salida detectada.
Si se observa hasta \(T\), el registro ideal de tiempo y censura es
\begin{equation}
 \widetilde\tau_D=\min\{\tau_D,T\},\qquad
 \Delta_D=\mathbf 1_{\{\tau_D\leq T\}}.
 \label{eq:reliability-censored-exit}
\end{equation}
En el cálculo, estos campos se estiman con un detector de eventos cuya
función de cruce, dirección, tolerancia e interpolación deben declararse. La
ausencia de salida detectada se registra como censura derecha. La afirmación
\(\tau_D>T\) sobre la órbita exacta requiere controlar todos los posibles
cruces entre muestras; sin ese control sólo se afirma que el detector no
encontró salida a su resolución. Un fallo del integrador se distingue de una
trayectoria observada sin escape.

Para precisar la detección, supóngase localmente
\(D=\{x:g(x)\leq0\}\), con \(g\in C^1\). Una salida transversal satisface
\[
 g(x(\tau_D))=0,\qquad
 \frac{\dd}{\dd t}g(x(t))\bigg|_{t=\tau_D}
 =\nabla g(x(\tau_D))\cdot f(x(\tau_D))>0.
\]
Si dos muestras consecutivas cumplen \(g_j\leq0<g_{j+1}\), la interpolación
lineal \(g_j+(t-t_j)(g_{j+1}-g_j)/(t_{j+1}-t_j)\) tiene el cero
\[
 \widehat\tau_D=t_j-
 \frac{g_j}{g_{j+1}-g_j}(t_{j+1}-t_j).
\]
Esta fórmula aproxima el evento de la curva interpolada. Si el cero exacto
es único en ese intervalo, la derivada de \(g(x(t))\) está acotada
inferiormente por \(\nu>0\), y el error uniforme de la función de evento es
\(\varepsilon_g\), el teorema del valor medio da un error de raíz a lo sumo
\(\varepsilon_g/\nu\), más la tolerancia temporal del buscador de ceros.
Aquí \(\varepsilon_g\) debe incluir el error de trayectoria y el de
interpolación. Sin transversalidad o sin control entre muestras pueden
omitirse salidas y reentradas, incluso si los signos en los nodos coinciden.

Para un conjunto medible de datos \(U\subset D\) y una medida declarada
\(\mu\) con \(0<\mu(U)<\infty\), la función de supervivencia es
\begin{equation}
 S_D(t)=
 \frac{\mu\{x_0\in U:\tau_D(x_0)>t\}}{\mu(U)}.
 \label{eq:reliability-survival}
\end{equation}
Como \(\{\tau_D>t_2\}\subseteq\{\tau_D>t_1\}\) si \(t_2>t_1\),
\(S_D\) es no creciente y pertenece a \([0,1]\). La finitud de \(\mu(U)\)
y la continuidad de la medida sobre conjuntos decrecientes dan
\[
 \lim_{t\to\infty}S_D(t)
 =\frac{\mu\{x_0\in U:\tau_D(x_0)=\infty\}}{\mu(U)}.
\]
Este límite asintótico no se determina a partir de un único horizonte finito.

Para un banco de \(N\) semillas con pesos iguales, la medida empírica es
\(\mu_N=N^{-1}\sum_{j=1}^N\delta_{x_0^{(j)}}\). Con un horizonte común
\(T\), su supervivencia observada se calcula, para \(0\leq t<T\), mediante
\begin{equation}
 \widehat S_D(t)=\frac1N\sum_{j=1}^N
                      \mathbf1_{\{\widetilde\tau_D^{(j)}>t\}}.
 \label{eq:reliability-empirical-survival}
\end{equation}
En efecto, \(\min(\tau_D^{(j)},T)>t\) equivale a
\(\tau_D^{(j)}>t\) cuando \(t<T\). El indicador de censura permite
separar salidas detectadas de observaciones detenidas en \(T\); en este
último tiempo ya no se usa la identidad anterior para inferir supervivencia
posterior. Con horizontes distintos se necesita un tratamiento estadístico de
la censura y sus hipótesis. Los intervalos de confianza requieren un modelo
de muestreo; una malla determinista proporciona conteos descriptivos a una
resolución fija.
En regímenes hiperbólicos puede aparecer
\(S_D(t)\sim C\e^{-\kappa t}\), donde \(\kappa\) es una tasa de escape; en
regímenes no hiperbólicos son posibles colas no exponenciales. Ajustar una
recta en una sola escala no demuestra ninguna de las dos leyes. Deben variarse
la región \(D\), la distribución de semillas, el paso y la ventana de ajuste.

En Rabinovich--Fabrikant se han descrito trayectorias ocultas con
comportamiento caótico transitorio \cite{Danca2016HiddenTransient}.
La denominación \emph{transitorio oculto} distingue estos conjuntos de los
atractores cuya atracción asintótica ha sido establecida
\cite{GuanXie2025}.

\subsection{Crisis y cambio del destino asintótico}

Una crisis es una bifurcación global en la que un conjunto caótico entra en
contacto con una silla u órbita periódica inestable y, en particular para una
crisis de frontera, con la variedad estable que forma parte de la frontera de
su cuenca. Las tres configuraciones que deben distinguirse son:
\begin{description}[style=nextline,leftmargin=3.5em]
 \item[Crisis interior] el atractor cambia bruscamente de tamaño al acceder a
 nuevas regiones; antes y después sigue existiendo un atractor.
 \item[Crisis de frontera] el atractor colisiona con la frontera de su cuenca
 y deja de atraer; la dinámica caótica previa puede sobrevivir como silla y
 producir transitorios largos.
 \item[Crisis de fusión] dos o más componentes atractoras se unen al variar un
 parámetro.
\end{description}
Una nube más grande o una caída abrupta del tiempo de residencia sólo sugiere
una crisis. Para sostener el mecanismo se necesita localizar el objeto
inestable implicado y seguir su intersección con el conjunto o con la frontera
\cite{Ott2002Chaos,TelGruiz2006Chaotic}.

En un problema Caputo, \(D\), la silla y sus conjuntos estable e inestable
deben formularse para una semidinámica bien definida en un espacio admisible
invariante \(X_{\mathrm{adm}}\subset\mathfrak C\), con la topología relativa
declarada en \cref{sec:geometria-topologia-caputo}. La realización en todo
\(\mathfrak C\) de \cref{eq:geofo-history-space} no basta para atribuirle
atracción local. Un escape de la proyección en
\(\mathbb R^n\) puede ser un observable útil, pero dos perfiles con igual
estado terminal pueden tener tiempos de salida distintos. Toda curva de
supervivencia fraccionaria debe fijar la familia de historias muestreadas y
separar continuación con historia de reinicio.

\subsection{Protocolo CT0--CT6 para distinguir atracción y transitorios prolongados}
\label{subsec:protocolo-transitorios}

\begin{description}[style=nextline,leftmargin=3.4em,labelwidth=2.9em]
 \item[CT0] Fijar campo, parámetros, operador, terminal inferior,
  inicialización, región de observación, detector de escape, tolerancias y
  tratamiento de la censura.
 \item[CT1] Verificar convergencia de la trayectoria o de los observables con
 \(h,h/2\) y, en una réplica corta, \(h/4\); usar un segundo integrador que
 represente el mismo PVI.
 \item[CT2] Calcular FTLE en varias longitudes y posiciones de ventana;
 conservar la distribución y no sólo la media.
 \item[CT3] Extender el horizonte de cada destino irregular y medir
  los tiempos e indicadores de \cref{eq:reliability-censored-exit}. Una
  trayectoria sin salida detectada se registra como censurada, no como prueba
  de \(\tau_D=\infty\); se controla la resolución de los cruces.
 \item[CT4] Muestrear un conjunto de semillas y estimar \(S_D(t)\) con
  el esquema de censura y, si el muestreo lo justifica, intervalos de confianza;
  repetir al cambiar \(D\) y el refinamiento. Para medidas empíricas de mallas
  se reportan los conteos y su resolución.
 \item[CT5] Para clasificar el conjunto como atractor, comprobar invariancia
 numérica de la cola
  y convergencia desde perturbaciones muestreadas en la familia de datos
  declarada, separando regímenes coexistentes. La atracción de un abierto de
  estados o de una vecindad relativa de \(X_{\mathrm{adm}}\) requiere una
  justificación adicional a la muestra finita.
 \item[CT6] Sólo después aplicar el sondeo de vecindades de equilibrios. La
 ocultedad es una afirmación de cuenca, no una consecuencia del tiempo de
 residencia, del FTLE o de una crisis.
\end{description}

\begin{longtable}{@{}L{0.22\linewidth}L{0.31\linewidth}L{0.37\linewidth}@{}}
\caption{Niveles de evidencia dinámica y controles mínimos.}
\label{tab:reliability-promotion}\\
\toprule
Afirmación & Evidencia mínima & Lo que todavía no demuestra\\
\midrule
\endfirsthead
\toprule
Afirmación & Evidencia mínima & Lo que todavía no demuestra\\
\midrule
\endhead
Trayectoria del PVI & Refinamiento, misma especificación e integrador independiente
o defecto controlado & Seguimiento punto a punto para tiempo arbitrario.\\
Expansión finita & Distribución FTLE robusta en ventanas declaradas &
Exponente asintótico, entropía o atracción.\\
Caos transitorio & Irregularidad robusta seguida de escapes y supervivencia
cuantificada & Existencia o hiperbolicidad de una silla sin análisis
geométrico.\\
Atractor numérico & Cola aparentemente invariante y atracción robusta desde
un abierto muestreado & Prueba matemática global de invariancia o
atracción.\\
Atractor caótico & Atracción más diagnósticos caóticos convergentes &
Herradura o entropía topológica sin argumento adicional.\\
Atractor oculto & Atractor identificado y ausencia de contacto en sondeos
anidados de todos los equilibrios & Separación global de cuencas a partir de
un conjunto finito de sondas.\\
\bottomrule
\end{longtable}

Las clases «otro destino» y «transitorio no resuelto» son resultados
científicamente válidos y se conservan cuando la evidencia no permite asignar
un comportamiento regular o caótico. En particular, las semillas FPP--A de
\cref{sec:perpetual-points} pueden entrar cerca de una silla caótica o salir
de ella; la cancelación local del término \(h^{2q}\) no decide el destino
asintótico.

\subsection{Índice de Conley y certificación mediante discretizaciones}

El índice de Conley permite estudiar conjuntos invariantes aislados sin
conocer todas sus órbitas. Sea \(N\) una vecindad compacta. Es aislante si
\begin{equation}
 \operatorname{Inv}(N)\subset\operatorname{int}N.
 \label{eq:reliability-isolating-neighborhood}
\end{equation}
Un par índice \((N_1,N_0)\) separa una región que retiene la dinámica de su
conjunto de salida. El índice de Conley es, esquemáticamente, el tipo de
homotopía apuntado del cociente \(N_1/N_0\). Un índice no trivial implica que
el bloque contiene un conjunto invariante; la invariancia por continuación
permite transportarlo mientras se conserve el aislamiento
\cite{Jost2005Dynamical,MischaikowEtAl2016Conley}.

La aplicación del índice requiere un bloque compacto con condición de
salida verificada; un flujo, semiflujo o mapa de retorno bien definido; un par
índice válido; y la relación entre el conjunto aislado y la cuenca objetivo.
En Caputo, el bloque pertenece a un espacio funcional admisible de memoria.
Si se obtiene de una elevación finita, su transporte al problema original
requiere una homotopía en un espacio común, admisibilidad y aislamiento
uniforme, según \cref{camp:caputo-lift}. La certificación computacional
combina este argumento con los métodos de cajas de
\cref{sec:localizacion-geometrico-topologica} y las propiedades del sistema
estudiado. Los controles finitos de paso, modos y horizonte aportan evidencia
numérica, mientras la atracción, el caos y la ocultedad requieren sus
respectivas demostraciones.

\subsection{Cuestiones de síntesis sobre confiabilidad dinámica}

\begin{enumerate}
 \item ¿Por qué un exponente positivo en tiempo finito no prueba caos
 asintótico?
 \item ¿Qué diferencia hay entre error pequeño, sombreado y conservación del
 destino de cuenca?
 \item ¿Cómo puede una silla caótica producir una nube que parece atractora
 durante miles de iteraciones?
 \item ¿Qué cambia al definir escape y supervivencia para Caputo?
 \item ¿Qué observación distingue una crisis de frontera de una expansión
 interior del atractor?
 \item ¿Qué hipótesis faltan entre un grafo de cajas y un índice de Conley
 válido?
\end{enumerate}

La resolución de estas cuestiones requiere citar la definición, indicar el
espacio de fases correspondiente, identificar el observable finito y delimitar
la afirmación que ese observable todavía no autoriza.
